\section{The multiplication ideal and intrinsic torsion}
\label{sec:fitting}
We first calculate on the whole Grassmannian. In this section
$R\in G_e^\circ$ is fixed, and the subscript $R$ is omitted.

\begin{proposition}[The full Fitting presentation]\label{prop:full-fitting}
On a frame chart write $\mathcal K\to V\oplus S$ as
$\binom{M}{A_0}$. For every $m\ge2$,
\begin{equation}\label{eq:full-matrix-ideal}
 \II_{\widehat D}=(\det M)I_p(\gamma\Sym^2M).
\end{equation}
The reduction of $\widehat D$ is the integral divisor
$\widehat\Delta$. Its smooth locus is the rank-$e-1$ stratum.
Consequently $D$ is the restriction of $\widehat D$ to the smooth
locus of its reduction, and this restriction is intrinsic.
\end{proposition}
\begin{proof}
Cube-zero multiplication and the unit show, over the coefficient ring
and not only on fibres, that the image of every degree $m\ge2$ is
$\C1+\mathcal K+\mathcal K^2$. Its presentation is
\[
 \begin{pmatrix}
  1&0&0\\0&M&0\\0&A_0&\gamma\Sym^2M
 \end{pmatrix}.
\]
Every nonzero maximal minor uses the scalar column and all $e$
columns meeting the $V$-rows. Expanding on these rows proves
\eqref{eq:full-matrix-ideal}. The graph block $A_0$ has no effect on
this ideal. Off $\det M=0$, the quadratic block is surjective because
$\gamma$ is surjective, so the ideal is the unit ideal. On the divisor
it vanishes. The divisor is the usual irreducible Schubert divisor:
the rank-$e-1$ stratum is dense and integral, and every lower-rank
plane is a specialization of that stratum. Equivalently the standard
Schubert charts give its reduced determinantal equations.

Near a point with projection rank $r$, an invertible $r$-minor
reduces the projection to $\diag(I_r,T)$, where the entries of the
$(e-r)\times(e-r)$ block $T$ are independent transverse coordinates.
One can see their independence by varying the $V$-components of a
basis of $K\cap S$ while keeping an invertible $S$-frame block fixed.
The determinant equation is $\det T$. Its gradient is nonzero
exactly at residual corank one. This proves the smooth-locus claim,
including its assertion on the Grassmannian rather than on a chosen
frame space.
\end{proof}

On $\Delta$ put $H=\im(K\to V)$ and $\Lambda=K\cap S$.
These give a morphism to $P\times\Pj(S)$. Define
$E=\Delta\times_PY$. The graph of $K/\Lambda$ yields isomorphisms
\begin{align}
 \Delta&\simeq\Tot_{P\times\Pj(S)}\Hom(\mathcal H,S/\Lambda),
       \label{eq:graph-delta}\\
 E&\simeq\Tot_{Y\times\Pj(S)}\Hom(\mathcal H,S/\Lambda).
       \label{eq:graph-E}
\end{align}
Indeed the inverse construction takes the inverse image of a graph
$H\to S/\Lambda$ in $H\oplus S$. Both constructions commute with
base change. In particular $E$ and $\Delta$ are smooth integral,
of dimensions $ep-2$ and $ep-1$.

\begin{theorem}[Ramification normal form]\label{thm:normal-form}
On all of $U$ the ideal of $D$ is
\begin{equation}\label{eq:normal-global}
 \II_D=\II_\Delta\II_E=\II_\Delta\cap\II_E^2.
\end{equation}
The associated supports are exactly $\Delta$ and $E$. The nilradical
is the invertible $\OO_E$-module $\NN=\OO_E(-\Delta)$ and is
square-zero. Its annihilator recovers $E$ as a scheme.
\end{theorem}
\begin{proof}
Near projection corank one, Gaussian operations make
$M=\diag(I_{e-1},t)$. The quadratic block becomes
$[A,tB,t^2c]$, with $A|_{t=0}=\gamma|_{\Sym^2H}$.
Equation~\eqref{eq:firstorder} says $[A,B]$ is onto locally; choose
$AU+BV=I_p$. Multiplication by $t$ then shows that
$[A,tB,t^2c]$ and $[A,tI_p]$ have exactly the same column module.
Their zeroth Fitting ideals agree. Off $Y$, $A$ is invertible and
\eqref{eq:full-matrix-ideal} is $(t)$. On $Y$, the restriction has
rank $p-1$ by Proposition~\ref{prop:relative-open}. A Schur complement
gives a residual row with ideal $(t,f)$, where $f\bmod t$ is a local
ramification equation. Thus the full ideal is
\begin{equation}\label{eq:tf}
 (t^2,tf)=t(t,f).
\end{equation}
The pair $t,f$ is a regular parameter pair cutting out $E$.
The identity $(t)\cap(t,f)^2=t(t,f)$ proves
\eqref{eq:normal-global}. These are primary ideals because their
associated graded rings are polynomial rings over the integral smooth
supports. At the generic point of $E$, neither primary component can
be omitted. Their radicals are distinct, giving exactly the stated
associated points. Finally the nilradical is
$(t)/t(t,f)$; cancellation of $t$ gives its module identification and
annihilator $(t,f)/t(t,f)$. Both are intrinsic, so the calculation glues.
\end{proof}

\begin{lemma}[Generic torsion as a sheaf]\label{lem:generic-torsion}
Let $Z$ be noetherian with integral reduction and generic point $\eta$.
For a coherent module $\mathcal F$ define
\[
 \operatorname{tor}_{\eta}(\mathcal F)
   =\ker\bigl(\mathcal F\longrightarrow
       i_*\widetilde{\mathcal F_\eta}\bigr),
 \qquad i:\Spec\OO_{Z,\eta}\longrightarrow Z.
\]
This kernel is coherent and is the largest coherent submodule with
zero generic stalk. It is preserved by isomorphisms of schemes.
\end{lemma}
\begin{proof}
On an affine open meeting $\eta$ this is the localization kernel at
the prime of $\eta$. It is a submodule of a finite module over a
noetherian ring, hence finite. Localization of that kernel on smaller
affine opens is the corresponding kernel, so these modules glue.
Every submodule with zero generic stalk lies in the kernel, and the
kernel itself has zero generic stalk. The definition involves no
ambient embedding or chosen primary decomposition.
\end{proof}

